\chapter{What Is an FPGA}
\label{ch:fpga}

\section{A circuit you can rewire without a soldering iron}

Imagine you need to build a device that turns on a lamp
when \emph{both} buttons are pressed. You take a chip with AND gates,
solder some wires, and it works. Tomorrow the customer asks: ``Let the lamp
light when \emph{exactly one} button is pressed.'' You will have to desolder
the AND gate, put in an ``exclusive OR'' gate, and move the wires around.

Now imagine a chip that contains thousands of
``blank'' logic gates and flip-flops and a huge network of
wires with controllable switches at every crossing. Which
blanks become what, and which switches are closed, is set by a
special configuration file that is loaded into the chip over a cable
in a fraction of a second. This is an \term{FPGA}.

\begin{ideabox}
An \textbf{FPGA} (\textbf{F}ield-\textbf{P}rogrammable \textbf{G}ate
\textbf{A}rray) is a chip whose internal
digital circuit can be changed many times after the chip has been manufactured.
\end{ideabox}

It is important to understand the difference right away: an FPGA does \emph{not run a program}, as a
processor does. It \emph{becomes a circuit}. Once the configuration is loaded, the
chip contains real logic gates and real wires,
and signals travel along them at the same time, just as in a circuit on a breadboard.

\begin{figure}[H]
\centering
\begin{tikzpicture}
  % Variant 1
  \begin{scope}
    \node[font=\sffamily\bfseries\small] at (1.4,1.9) {Configuration 1};
    \node[gAND] (g1) at (1.4,0) {};
    \draw[wire] (g1.input 1) -- ++(-0.8,0) node[left] {$a$};
    \draw[wire] (g1.input 2) -- ++(-0.8,0) node[left] {$b$};
    \draw[wire] (g1.output) -- ++(0.8,0) node[right] {$y=a\cdot b$};
  \end{scope}
  % Arrow
  \draw[-{Stealth[length=4mm]}, line width=2pt, FpgaOrange] (4.6,0) -- (7.0,0)
    node[midway, above=2mm, font=\sffamily\footnotesize, text=FpgaDark, align=center] {new configuration\\file};
  % Variant 2
  \begin{scope}[xshift=8.6cm]
    \node[font=\sffamily\bfseries\small] at (0.3,1.9) {Configuration 2};
    \node[gXOR] (g2) at (0.3,0) {};
    \draw[wire] (g2.input 1) -- ++(-0.8,0) node[left] {$a$};
    \draw[wire] (g2.input 2) -- ++(-0.8,0) node[left] {$b$};
    \draw[wire] (g2.output) -- ++(0.8,0) node[right] {$y=a\oplus b$};
  \end{scope}
  \draw[dashed, FpgaGray, rounded corners] (-1.0,-1.1) rectangle (13.0,2.4);
  \node[lbl, anchor=north] at (6.0,-1.2) {the same FPGA chip, the same pins};
\end{tikzpicture}
\caption{The same FPGA after loading different configurations}
\label{fig:reconfig}
\end{figure}

\begin{diybox}[One chip, two circuits]
\textbf{Task.} Repeat the experiment of Fig.~\ref{fig:reconfig}: first build an AND circuit, then ``rewire'' it into an ``exclusive OR'' without changing the inputs and output.\\
\textbf{In Logra.} ``New project'' $\to$ ``From logic elements''; the \emph{Designer} tab; then \emph{Waveforms}.\\
\textbf{Expected result.} For AND, the output is 1 only when $a=b=1$; after replacing the gate with XOR, it is 1 when $a\ne b$.\\
\textbf{How to check.} On the \emph{Waveforms} tab, press ``Sweep inputs'' in both versions and compare the four combinations; on the RTL tab, make sure that only one element has changed.
\end{diybox}

\section{A little history}

The idea of programmable logic appeared in the 1970s. At first these were
small AND--OR arrays (PAL, PLA) in which fuses were blown.
Then came CPLDs (complex programmable logic devices), and in 1985 Xilinx
released the first FPGA, the XC2064, which had only 64 logic blocks.
Modern FPGAs contain millions of logic cells.

\begin{figure}[H]
\centering
\begin{tikzpicture}
  \draw[line width=2pt, FpgaBlue, -{Stealth[length=4mm]}] (0,0) -- (15,0);
  \foreach \x/\y/\t/\d in {
      0.8/1/1975/PLA, 3.2/-1/1978/PAL,
      5.6/1/1984/CPLD (Altera), 8/-1/1985/{Xilinx XC2064,\\ the first FPGA},
      10.4/1/2000s/{FPGAs with DSP\\and memory}, 13/-1/2020s/{Tang Nano 20K:\\ 20\,736 LUTs for pennies}} {
    \fill[FpgaOrange] (\x,0) circle (1.1mm);
    \draw[FpgaGray] (\x,0) -- (\x,\y*0.55);
    \node[font=\sffamily\bfseries\small, text=FpgaDark] at (\x,\y*0.85) {\t};
    \node[font=\sffamily\footnotesize, align=center, text=FpgaDark, anchor={\ifnum\y>0 south\else north\fi}] at (\x,\y*1.1) {\d};
  }
\end{tikzpicture}
\caption{The evolution of programmable logic}
\end{figure}

\begin{diybox}[The scale of a chip]
\textbf{Task.} Compare the 64 logic blocks of the first FPGA (1985) with the Tang Nano 20K chip.\\
\textbf{In Logra.} Open the ``Running light'' project and press ``Synthesize''; in the ``Real board'' group, open the \emph{Chip} tab, then ``Place and build'' and the ``Whole chip'' button.\\
\textbf{Expected result.} A map of the whole chip; the caption under it gives the size of the tile grid (X from 0 to 55, Y from 0 to 54), and above it are counters of the used LUT, FF, ALU, and IO resources.\\
\textbf{How to check.} Write down the number of used LUTs (15) and the total number of LUTs (20\,736, see Chapter~\ref{ch:inside}), and compute the percentage; compare it with the 64 blocks of the XC2064.
\end{diybox}

\section{FPGA, microcontroller, and custom chip}

To understand why FPGAs are needed, let us compare three ways of making a ``smart''
digital device.

\begin{description}
  \item[Microcontroller (MCU)] --- a ready-made processor that
    runs a program \emph{sequentially}, one instruction after another. For example,
    the autopilot of the Geoscan Pioneer quadcopter runs on an STM32 MCU.
  \item[Custom chip (ASIC)] --- a circuit ``burned'' into silicon at a
    factory. It is faster and more economical than anything else, but developing it costs
    millions of dollars, and it can no longer be changed.
  \item[FPGA] --- the golden mean: a hardware circuit, like an ASIC,
    that can be changed like the program of an MCU.
\end{description}

\begin{figure}[H]
\centering
\begin{tikzpicture}
  % MCU
  \begin{scope}
    \node[font=\sffamily\bfseries] at (2.6,3.6) {Microcontroller};
    \node[hblock, minimum width=16mm] (cpu) at (2.6,2.4) {CPU};
    \foreach \i/\t in {0/{$a\cdot b$}, 1/{$c+d$}, 2/{$e\oplus f$}, 3/{$g\cdot h$}} {
      \node[block, minimum width=14mm, minimum height=6mm, font=\sffamily\footnotesize]
        (t\i) at (2.6, 1.2-\i*0.8) {\t};
    }
    \draw[arr] (cpu) -- (t0);
    \draw[arr] (t0) -- (t1); \draw[arr] (t1) -- (t2); \draw[arr] (t2) -- (t3);
    \draw[arr, FpgaOrange] (0.9,1.3) -- (0.9,-1.3) node[midway, left, font=\sffamily\footnotesize, align=right] {time\\4 steps};
  \end{scope}
  % FPGA
  \begin{scope}[xshift=7.2cm]
    \node[font=\sffamily\bfseries] at (2.8,3.6) {FPGA};
    \foreach \i/\t in {0/{$a\cdot b$}, 1/{$c+d$}, 2/{$e\oplus f$}, 3/{$g\cdot h$}} {
      \node[gblock, minimum width=12mm, minimum height=6mm, font=\sffamily\footnotesize]
        (p\i) at (\i*1.5+0.5, 1.2) {\t};
      \draw[arr] (\i*1.5+0.5, 2.5) -- (p\i);
      \draw[arr] (p\i) -- (\i*1.5+0.5, 0);
    }
    \draw[arr, FpgaOrange] (-0.5,1.5) -- (-0.5,0.9) node[midway, left, font=\sffamily\footnotesize, align=right] {1 step};
    \node[lbl, align=center] at (2.75,-0.6) {all four operations are\\ separate circuits that work \emph{at the same time}};
  \end{scope}
\end{tikzpicture}
\caption{A processor performs actions one after another; an FPGA performs them in parallel}
\label{fig:par}
\end{figure}

Figure~\ref{fig:par} shows the main superpower of FPGAs: \term{parallelism}.
If you need to perform 1000 independent operations, a processor needs
1000 steps, while in an FPGA you can build 1000 circuits that finish in a single
clock cycle.

\begin{figure}[H]
\centering
\begin{tikzpicture}
\begin{axis}[
  width=0.82\textwidth, height=6.2cm,
  xlabel={Number of independent operations}, ylabel={Time, $\mu$s},
  xmin=0, xmax=1000, ymin=0, ymax=12,
  grid=major, grid style={FpgaGray!25},
  legend pos=north west, legend style={font=\sffamily\small, draw=none, fill=white, fill opacity=0.8},
  tick label style={font=\small}, label style={font=\sffamily\small},
]
  \addplot[very thick, FpgaRed, domain=0:1000, samples=2] {x*0.0104};
  \addlegendentry{Microcontroller, 96\,MHz (1 operation per cycle)}
  \addplot[very thick, FpgaGreen, domain=0:1000, samples=2] {0.037};
  \addlegendentry{FPGA, 27\,MHz (all operations in 1 cycle)}
\end{axis}
\end{tikzpicture}
\caption{Processing time as the number of operations grows (a simplified estimate)}
\end{figure}

\begin{table}[H]
\centering
\caption{Comparison of the three approaches}
\small
\begin{tabularx}{\textwidth}{>{\bfseries}X c c c}
\toprule
 & \textbf{Microcontroller} & \textbf{FPGA} & \textbf{ASIC} \\
\midrule
How it works          & a program, step by step & a circuit, in parallel & a circuit, in parallel \\
Response time         & microseconds        & nanoseconds        & nanoseconds \\
Can be changed?       & yes                 & yes                & no \\
Cost of one unit      & low                 & medium             & enormous \\
Development effort    & low                 & medium             & very high \\
Power consumption     & low                 & medium             & lowest \\
\bottomrule
\end{tabularx}
\end{table}

\begin{diybox}[Parallelism with your own eyes]
\textbf{Task.} Make sure that independent circuits in an FPGA work at the same time.\\
\textbf{In Logra.} \emph{Designer}: put two independent elements, AND and XOR, in one project, with shared inputs $a$ and $b$ and two different outputs; then open the \emph{Waveforms} tab.\\
\textbf{Expected result.} Both outputs change at the same moment, as in Fig.~\ref{fig:par}; there is no ordering at all.\\
\textbf{How to check.} Use ``Sweep inputs'' and, on the waveform, find a combination where the outputs differ and combinations where they match. \stars{2} Add a third element and make sure that the response time has not grown.
\end{diybox}

\section{Where FPGAs are used}

FPGAs are chosen where one of three things is needed: \emph{very high speed},
\emph{a great deal of parallel work}, or \emph{exact timing}.

\begin{figure}[H]
\centering
\begin{tikzpicture}[
  app/.style={block, minimum width=48mm, minimum height=14mm, font=\sffamily\footnotesize},
]
  \node[hblock, circle, minimum size=24mm, font=\sffamily\bfseries] (c) at (0,0) {FPGA};
  \foreach \a/\t/\s in {
     90/{Video processing\\and machine vision}/app,
     30/{Communications: 5G, radar,\\software-defined radio}/app,
    -30/{Motor control,\\PWM, drone sensors}/app,
    -90/{Prototyping\\processors (ASICs)}/app,
   -150/{Interface\\bridging}/app,
    150/{AI and computing\\acceleration}/app} {
    \node[\s] (n\a) at (\a:4.1) {\t};
    \draw[thick, FpgaBlue!60] (c) -- (n\a);
  }
\end{tikzpicture}
\caption{The main application areas of FPGAs}
\end{figure}

\begin{itemize}
  \item \textbf{Digital signal processing}: filters, the Fourier
        transform, processing of audio and radio signals in real time.
  \item \textbf{Video}: generating the image for HDMI, processing frames from cameras,
        finding objects right ``on the fly.''
  \item \textbf{Communications and interfaces}: an FPGA can ``speak'' any
        protocol (UART, SPI, I$^2$C, Ethernet, USB), even one you invented yourself.
        Here a protocol is a set of rules by which two chips
        exchange bits over a few wires. We will study the simple serial
        protocols UART, SPI, and I$^2$C in Chapter~\ref{ch:uart}.
  \item \textbf{Robotics and UAVs}: precise generation of dozens of PWM signals (pulse-width modulation:
        fast pulses whose ``one'' fraction sets, for example, the speed
        of a motor; Chapter~\ref{ch:examples}),
        simultaneous polling of many sensors, video stabilization.
  \item \textbf{Prototyping}: before ordering an ASIC for
        millions of dollars, engineers test it on an FPGA.
  \item \textbf{Retro computers}: FPGAs are used to recreate old game
        consoles and computers (the MiSTer project), accurate down to the clock cycle.
\end{itemize}

\begin{ideabox}
An FPGA is ``modeling clay'' for digital circuit design. You describe a circuit,
a synthesis program lays it out on the resources of the chip, and within
a minute your circuit is already running in hardware.
\end{ideabox}

\begin{diybox}[Where an FPGA will go]
\textbf{Task.} Pick one example from the Logra catalog for each of three application areas from the figure.\\
\textbf{In Logra.} ``New project'' $\to$ ``From an example'' (the catalog has 275 projects): browse the categories and choose three examples, for instance one on PWM, one on communications, and one on displays.\\
\textbf{Expected result.} A table of three rows: the application area, the name of the example, and one sentence about what it does.\\
\textbf{How to check.} Open each example, press ``Synthesize'', and make sure that the RTL is built without errors.
\end{diybox}

\section{What Next}

We have found out that an FPGA becomes a circuit, and we know where such
circuits are used. A practical question remains: how do we control this
circuit? Designing needs a tool that describes the circuit, checks it, and
shows what it turned into. The next chapter, Chapter~\ref{ch:logra}, introduces
Logra: you will create a first project, walk through the tabs, and see the whole
path from idea to board. After that we start to work out what such circuits are
made of, beginning with the simplest thing, the signal (Chapter~\ref{ch:signals}).
